\section{Interdisciplinary Significance}
\label{sec:interdisciplinary}

The \Apeiron{} Framework is not confined to artificial intelligence; it has implications for several neighboring disciplines:

\begin{itemize}
    \item \textbf{Philosophy of Science.} The framework provides a computational instantiation of perspectival realism \cite{massimi2022} and epistemic iteration \cite{chang2012}. It offers a concrete model for how scientific knowledge could be constructed without human cognitive biases.
    \item \textbf{Complex Systems.} The seventeen layers constitute a formally specified hierarchy of emergence, from substrate to world-making. This provides a benchmark for theories of weak and strong emergence \cite{bedau1997}.
    \item \textbf{Category Theory.} The framework applies advanced categorical constructions (limits, colimits, adjunctions, monads) to a real-world architectural problem, demonstrating the practical utility of applied category theory \cite{fong2019}.
    \item \textbf{Ethics and Governance.} The embedded ethical framework (Layers 15--17) offers a novel approach to AI alignment: making reversibility, auditability, and participatory governance first-class architectural principles rather than external constraints.
    \item \textbf{Synthetic Biology and Programmable Matter.} Layer~15 provides a formal vocabulary for describing recursive morphogenesis, with potential applications in the design of self-replicating systems and the governance of synthetic biology.
\end{itemize}